\setminted{breaklines=true, breakanywhere=true, fontsize=\footnotesize, fontfamily=lmtt,
xleftmargin=0pt, xrightmargin=0pt, framesep=1pt}
\begin{figure*}[ht]
  \begin{subfigure}[b]{0.36\textwidth}
    \centering
    \begin{minted}{python}
for bx in range(N0, "blockIdx.x"):
  for i1 in range(N1):
    xmax[bx*N1+i1] = -inf
  for j0 in range(M0):
    for i1, j1, k in grid(N1, M1, K):
      i, j = bx*N1+i1, j0*M1+j1
      mm[i, j] += Q[i, k] * K[j, k]
    for i1, j1 in grid(N1, M1):
      i, j = bx*N1+i1, j0*M1+j1
      xmax[i] = max(
        xmax[i], mm[i, j])
    \end{minted}
    \caption{Scalar IR, with scalar memory accesses and computation under loop nests.
      This IR facilitates high-level optimizations such as operator fusion,
    which are central to \NAME{}'s auto-scheduler.}
  \end{subfigure}
  \hspace{2pt}
  \begin{subfigure}[b]{0.32\textwidth}
    \centering
    \begin{minted}{python}
bx = block("blockIdx.x", 0, N0)
vmax_init = tile(-inf, [N1])
vmax = for j0 in loop(M0)
    params (vmax_1 = vmax_init) {
  il, ir = bx*N1, (bx+1)*N1
  jl, jr = j0*M1, (j0+1)*M1
  vmm = dot(Q[il:ir, 0:K],
    K[jl:jr, 0:K].T)
  vmax_local = max(vmm, axis=1)
  vmax = maximum(
    vmax_local, vmax_1)
  yield vmax
}
store(vmax, xmax[bx*N1:bx*N1+N1])
    \end{minted}
    \caption{Virtual register (VR) tile IR facilitates expression rewrites.
    Tiles are expressed as values instead of tensor accesses.}
    \label{fig:vr-ir-example}
  \end{subfigure}
  \hspace{2pt}
  \begin{subfigure}[b]{0.31\textwidth}
    \centering
    \begin{minted}{python}
bx = block("blockIdx.x", 0, N0)
for j0 in range(M0):
  il, ir = bx*N1, (bx+1)*N1
  jl, jr = j0*M1, (j0+1)*M1
  xmm[0:N1, 0:M1] = dot(
    Q[il:ir, 0:K],
    K[jl:jr, 0:K].T)
  xmax_local = tile(-inf, [N1])
  xmax_local[0:N1] = max(
    xmm[:, :], axis=1)
  xmax[0:N1] = maximum(
    xmax_local[:], xmax[:])
    \end{minted}
    \caption{Memory access (MA) tile IR desugars the VR-tile IR and exposes data
      movement between tensors. \NAME{} uses this IR to interface with its tile
      optimizers (Triton, Tawa, and TileLang).
    }
  \end{subfigure}
  \vspace*{-1.5em}
  \caption{An example program in \NAME{}'s three IRs.}
  \label{fig:3ir-example}
  \vspace*{-1em}
\end{figure*}
